% Einheit 1 von 4 – Terme aufstellen und berechnen (bank/terme/e1.jsonl, zusammenbau v0.1)
%% TODO Zweigzeile Teil 1: was hier gelernt wird (Satz in Schülersprache aus dem Einheitstitel) – Einheit 1
\einheitenkopf[e1]{Einheit 1 von 4 · Terme aufstellen und berechnen}
\zweigzeile{ab Kl. 6, je nach Buch bis Kl. 7 · P10 oft}
\setcounter{aufgabe}{6}

%% TODO Titel: Ich-kann-Satz fehlt in der Bank (Platzhalter: Kettenname) – Nr. 7
%% TODO Anweisung: ein Satz über der ersten Teilaufgabe fehlt in der Bank – Nr. 7
\begin{aufgabe}{Term aufstellen}
%% TODO \rechenplatz{n}: Schrittzahl des Grundfalls fehlt in der Bank (nur bei fester Schreibform, 2.3 b)
\begin{teile}
\teil Das Vierfache einer Zahl $x$, um 1 vermindert – welcher Term passt? \\ \kreuz{$4x - 1$} \\ \kreuz{$4 \cdot (x - 1)$} \\ \kreuz{$1 - 4x$}
\teil Das Dreifache einer Zahl $x$, um 8 vergrößert – welcher Term? \leerfeld
\teil Flächeninhalt des Rechtecks (Maße in cm) – welcher Term?

\viereck[seiten={5,x,5,x}]{(0,0)}{(5,0)}{(5,2.5)}{(0,2.5)}
\leerfeld
\teil Ein Heft kostet $a$ Euro, ein Stift $b$ Euro. Lena kauft 4 Hefte und 3 Stifte. Gesamtpreis in Euro – welcher Term? \leerfeld
\teil Umfang eines Rechtecks mit den Seiten $x$ cm und 9 cm – Term, zusammengefasst? \leerfeld
\teil Verdreifache eine Zahl $x$, addiere 2 und verdopple die Summe – welcher Term? \leerfeld
\teil Die Differenz aus dem Dreifachen einer Zahl $x$ und 7 wird verdoppelt. Welcher Term passt? \hfill (P10 2023 OS) \\ \kreuz{$(3x - 7) : 2$} \\ \kreuz{$3x : 7 \cdot 2$} \\ \kreuz{$2 \cdot (3x - 7)$} \\ \kreuz{$3x - 7 \cdot 2$}
\end{teile}
\end{aufgabe}

%% TODO Titel: Ich-kann-Satz fehlt in der Bank (Platzhalter: Kettenname) – Nr. 8
%% TODO Anweisung: ein Satz über der ersten Teilaufgabe fehlt in der Bank – Nr. 8
\begin{aufgabe}{Termwert berechnen}
\begin{teile}
\teil $4x - 3$ für $x = 5$ – Termwert? \leerfeld
\end{teile}
\end{aufgabe}

%% TODO Titel: Ich-kann-Satz fehlt in der Bank (Platzhalter: Kettenname) – Nr. 9
%% TODO Anweisung: ein Satz über der ersten Teilaufgabe fehlt in der Bank – Nr. 9
\begin{aufgabe}{Situation zu Term angeben}
\begin{teile}
\teil Erfinde eine Situation, zu der der Term $12 + 3x$ passt. Wofür steht $x$? Was bedeutet der Termwert für $x = 2$?
\end{teile}
\end{aufgabe}

%% TODO Titel: Ich-kann-Satz fehlt in der Bank (Platzhalter: Kettenname) – Nr. 10
%% TODO Anweisung: ein Satz über der ersten Teilaufgabe fehlt in der Bank – Nr. 10
\begin{aufgabe}{Term aufstellen – Fehler finden, Darstellungswechsel, Anwendung}
\begin{teile}
\teil Ben berechnet den Termwert von $4 + 2x$ für $x = 3$. Finde den Fehler und rechne richtig. \rechnung{4 + 2 \cdot 3 &= 6 \cdot 3 \\ &= 18}
\teil Welcher Term gehört zum Termbaum?

\termbaum{\cdot}{4}{\tb{+}{x}{6}}
\leerfeld
\teil Ein Stromtarif kostet 12 € Grundpreis im Monat und 0,30 € je Kilowattstunde. Stelle einen Term für die Monatskosten in Euro bei $n$ Kilowattstunden auf. Was kostet ein Monat mit 150 Kilowattstunden?
\end{teile}
\end{aufgabe}
